% !TeX program = xelatex
% ============================================================
% pz/pz.tex — Пояснительная записка (ГОСТ 19.404-79)
%
% Главный содержательный документ курсового проекта: описывает
% назначение, постановку задачи, архитектуру и функционирование
% программы, обоснование выбора средств и ожидаемые технико-
% экономические показатели. По составу заменяет ВКР в курсовом
% комплекте (та же роль «пояснительного» документа, но по ЕСПД
% ГОСТ 19.404-79, а не по ГОСТ 7.32).
%
% Жёлтые блоки \hseFill{…} — места, которые нужно заменить своим
% текстом. Данные (ФИО, название, коды) приходят из настроек
% проекта через макросы \hse… (common/meta.tex).
% ============================================================
\documentclass[12pt,a4paper]{extarticle}
\input{../common/preamble}
\newcommand{\docname}{Пояснительная записка}
% ЕСПД-код пояснительной записки — тип «81» (ГОСТ 19.103-77).
\newcommand{\doccode}{\hseDocCodeBase\ 81 01-1}
\newcommand{\doccodelu}{\hseDocCodeBase\ 81 01-1-ЛУ}
\input{../common/commands}
\input{../common/styles}
\begin{document}
\sloppy
\input{../common/title_template}


% ============================================
% АННОТАЦИЯ
% ============================================
\section*{АННОТАЦИЯ}
\addcontentsline{toc}{section}{АННОТАЦИЯ}

Пояснительная записка~--- документ, содержащий сведения об~используемых алгоритмах, принципах функционирования и~другой информации, касающейся работы программы.

Настоящая пояснительная записка «\hseProjectName» содержит следующие разделы: «Введение», «Назначение и~область применения», «Технические характеристики», «Ожидаемые технико-экономические показатели», «Источники, использованные при~разработке».

В разделе «Введение» указаны наименование программы и~условное обозначение темы разработки, а~также основание для разработки.

В разделе «Назначение и~область применения» приведены назначение программы и~краткая характеристика области применения программы.

Раздел «Технические характеристики» содержит следующие подразделы: постановка задачи на~разработку программы; описание алгоритма и~функционирования программы; описание и~обоснование выбора метода организации входных и~выходных данных; описание и~обоснование выбора состава технических и~программных средств.

В разделе «Ожидаемые технико-экономические показатели» указаны предполагаемая потребность и~экономические преимущества разработки по~сравнению с~отечественными и~зарубежными образцами или~аналогами.

\clearpage

% ============================================
% СПИСОК ИСПОЛЬЗУЕМЫХ ГОСТ
% ============================================
\noindent
Оформление программного документа «Пояснительная записка» произведено по~требованиям ЕСПД:

\begin{enumerate}[label=\arabic*), leftmargin=1.25cm]
\item ГОСТ 19.101-77 Виды программ и~программных документов\cite{gost19101}.
\item ГОСТ 19.103-77 Обозначения программ и~программных документов\cite{gost19103}.
\item ГОСТ 19.104-78 Основные надписи\cite{gost19104}.
\item ГОСТ 19.105-78 Общие требования к~программным документам\cite{gost19105}.
\item ГОСТ 19.106-78 Требования к~программным документам, выполненным печатным способом\cite{gost19106}.
\item ГОСТ 19.404-79 Пояснительная записка. Требования к~содержанию и~оформлению\cite{gost19404}.
\end{enumerate}

Изменения к~данной пояснительной записке оформляются согласно ГОСТ\,19.603-78\cite{gost19603}, ГОСТ\,19.604-78\cite{gost19604}.

\clearpage

% ============================================
% СОДЕРЖАНИЕ
% ============================================
\hypersetup{linkcolor=black}
\tableofcontents
\hypersetup{linkcolor=blue}
\thispagestyle{fancy}
\clearpage

% ============================================
% 1. ВВЕДЕНИЕ
% ============================================
\section{ВВЕДЕНИЕ}

\subsection{Наименование программы}

% Наименования подставляются автоматически: полное — из названия проекта,
% английское и краткое — из настроек проекта.
Наименование программы~--- «\projectname».

Наименование программы на английском языке~--- «\hseProjectNameEn».

\subsection{Условное обозначение темы разработки}

Условное обозначение темы разработки~--- «\hseProjectNameShort».

\subsection{Основание для разработки}

% Подсказка: сошлитесь на учебный план направления и на утверждённую
% академическим руководителем тему курсового проекта.
Основанием для разработки является учебный план подготовки бакалавров по~направлению \hseFill{код и наименование направления подготовки} и~утверждённая академическим руководителем образовательной программы «\hseSpecialization» тема курсового проекта.

\clearpage

% ============================================
% 2. НАЗНАЧЕНИЕ И ОБЛАСТЬ ПРИМЕНЕНИЯ
% ============================================
\section{НАЗНАЧЕНИЕ И ОБЛАСТЬ ПРИМЕНЕНИЯ}

\subsection{Назначение программы}

\subsubsection{Функциональное назначение}

% Подсказка: перечислите основные функции, которые выполняет программа.
Функциональным назначением программы является \hseFill{основное функциональное назначение программы}. Данная работа реализует \hseFill{перечислите ключевые компоненты и функции}.

\subsubsection{Эксплуатационное назначение}

% Подсказка: опишите, кем и в каких условиях программа эксплуатируется.
Данная программа может применяться для \hseFill{сценарии эксплуатации}. Целевой аудиторией проекта являются \hseFill{перечислите целевых пользователей}.

\subsection{Краткая характеристика области применения}

% Подсказка: 1–2 абзаца о том, для чего и в какой предметной области
% применяется программа, какие задачи она решает.
Программа предназначена для использования в~\hseFill{предметная область и системы, где применяется программа}. Она может применяться как \hseFill{сценарий 1}, так и \hseFill{сценарий 2}.

\clearpage

% ============================================
% 3. ТЕХНИЧЕСКИЕ ХАРАКТЕРИСТИКИ
% ============================================
\section{ТЕХНИЧЕСКИЕ ХАРАКТЕРИСТИКИ}

\subsection{Постановка задачи на разработку программы}

% Подсказка: сформулируйте решаемую задачу и приведите функциональные
% требования (можно сослаться на требования из ТЗ, префикс ТЗ-Ф-01…).
Данная работа разрабатывается для \hseFill{цель разработки программы}. Функциональные требования к~программе приведены в~техническом задании и~детализированы ниже.

\begin{hseExample}
Пример перечня решаемых задач:
\begin{itemize}[nosep, leftmargin=1.25cm]
    \item \hseFill{задача 1};
    \item \hseFill{задача 2};
    \item \hseFill{дополните перечень задач}.
\end{itemize}
\end{hseExample}

\subsection{Описание алгоритма и функционирования программы}

% Подсказка: это центральный раздел ПЗ. Опишите архитектуру, ключевые
% алгоритмы, схемы данных, взаимодействие компонентов. Приветствуются
% диаграммы (компонентов, последовательностей, классов) и псевдокод.
\subsubsection{Архитектура программы}

\hseFill{Опишите архитектуру: компоненты, их назначение и взаимодействие. Приведите диаграмму архитектуры.}

\subsubsection{Описание функционирования}

\hseFill{Опишите основные сценарии работы программы: как обрабатываются запросы, каковы ключевые алгоритмы и потоки данных.}

\subsection{Описание и обоснование выбора метода организации входных и выходных данных}

% Подсказка: форматы, протоколы, кодировки, схемы данных; почему выбраны
% именно они.
\hseFill{Опишите форматы и протоколы входных и выходных данных, схему хранения данных и обоснуйте выбор.}

\subsection{Описание и обоснование выбора состава технических и программных средств}

% Подсказка: языки, фреймворки, СУБД, инфраструктура; обоснование выбора.
\hseFill{Перечислите использованные языки, фреймворки, библиотеки, СУБД и инфраструктуру; обоснуйте выбор каждого средства.}

\clearpage

% ============================================
% 4. ОЖИДАЕМЫЕ ТЕХНИКО-ЭКОНОМИЧЕСКИЕ ПОКАЗАТЕЛИ
% ============================================
\section{ОЖИДАЕМЫЕ ТЕХНИКО-ЭКОНОМИЧЕСКИЕ ПОКАЗАТЕЛИ}

\subsection{Предполагаемая потребность}

% Подсказка: опишите потенциальных пользователей и предполагаемую потребность.
Разрабатываемая программа может представлять интерес для \hseFill{целевая аудитория}. Потенциальными пользователями являются \hseFill{перечислите потенциальных пользователей}.

\subsection{Экономические преимущества разработки по~сравнению с~отечественными и~зарубежными образцами или~аналогами}

% Подсказка: кратко опишите преимущества разработки перед аналогами
% (можно с таблицей сравнения).
\hseFill{Опишите экономические и функциональные преимущества разработки по сравнению с существующими аналогами.}

\clearpage
% Страницы источников печатаются без нижнего штампа (как в реальных
% комплектах) — только номер листа и код документа сверху.
\pagestyle{appendix}

% ============================================
% ИСТОЧНИКИ, ИСПОЛЬЗОВАННЫЕ ПРИ РАЗРАБОТКЕ
% ============================================
\section*{ИСТОЧНИКИ, ИСПОЛЬЗОВАННЫЕ ПРИ РАЗРАБОТКЕ}
\addcontentsline{toc}{section}{ИСТОЧНИКИ, ИСПОЛЬЗОВАННЫЕ ПРИ РАЗРАБОТКЕ}

% Подсказка: добавьте ПЕРЕД ГОСТами источники по технологиям вашего
% проекта в алфавитном порядке, по образцу:
% \bibitem{fastapi}
% FastAPI [Электронный ресурс]. Режим доступа: \url{https://fastapi.tiangolo.com}, свободный. (дата обращения: 14.03.2026).
\renewcommand{\refname}{}
\begin{thebibliography}{99}
\bibitem{gost19101}
ГОСТ 19.101-77: Виды программ и~программных документов. // Единая система программной документации.~--- М.: ИПК Издательство стандартов, 2001.

\bibitem{gost19103}
ГОСТ 19.103-77: Обозначения программ и~программных документов. // Единая система программной документации.~--- М.: ИПК Издательство стандартов, 2001.

\bibitem{gost19104}
ГОСТ 19.104-78: Основные надписи. // Единая система программной документации.~--- М.: ИПК Издательство стандартов, 2001.

\bibitem{gost19105}
ГОСТ 19.105-78: Общие требования к~программным документам. // Единая система программной документации.~--- М.: ИПК Издательство стандартов, 2001.

\bibitem{gost19106}
ГОСТ 19.106-78: Требования к~программным документам, выполненным печатным способом. // Единая система программной документации.~--- М.: ИПК Издательство стандартов, 2001.

\bibitem{gost19404}
ГОСТ 19.404-79: Пояснительная записка. Требования к~содержанию и~оформлению. // Единая система программной документации.~--- М.: ИПК Издательство стандартов, 2001.

\bibitem{gost19603}
ГОСТ 19.603-78: Общие правила внесения изменений. // Единая система программной документации.~--- М.: ИПК Издательство стандартов, 2001.

\bibitem{gost19604}
ГОСТ 19.604-78: Правила внесения изменений в~программные документы, выполненные печатным способом. // Единая система программной документации.~--- М.: ИПК Издательство стандартов, 2001.
\end{thebibliography}

\clearpage

% ============================================
% ПЕРЕЧЕНЬ ТЕРМИНОВ И ОПРЕДЕЛЕНИЙ
% ============================================
\section*{ПЕРЕЧЕНЬ ТЕРМИНОВ И ОПРЕДЕЛЕНИЙ}
\addcontentsline{toc}{section}{ПЕРЕЧЕНЬ ТЕРМИНОВ И ОПРЕДЕЛЕНИЙ}

% Подсказка: добавьте термины предметной области и технологий вашего
% проекта. Порядок — алфавитный: сначала латиница, затем кириллица.
\begin{enumerate}[label=\arabic*., leftmargin=1.25cm]
\item \textbf{\hseFill{термин на латинице}}~--- \hseFill{определение термина}.
\item \textbf{\hseFill{ещё один термин}}~--- \hseFill{определение термина}.
\end{enumerate}

\clearpage


% ============================================
% ЛИСТ РЕГИСТРАЦИИ ИЗМЕНЕНИЙ
% ============================================
\input{../common/change_log}

\end{document}
